\section*{Hoofdstuk \ref{ch:g12:early-earth} --- De jonge Aarde en haar atmosfeer}

\begin{solution}{exo:g12:early-earth:1}
(a) \omterm{def:g12:early-earth:primitive-atmosphere}{ontgassing}; (b) een \omterm{def:g12:early-earth:stromatolite}{stromatoliet}; (c) een \omterm{def:g12:early-earth:bif}{gebande ijzerformatie}; (d) \omterm{def:g12:early-earth:red-beds}{rode
lagen}; (e) de \omterm{def:g12:early-earth:goe}{Grote Oxidatie}.
\end{solution}

\begin{solution}{exo:g12:early-earth:2}
Zirkoonkorrels die vloeibaar water aantonen (4.4 miljard jaar); de eerste \omterm{def:g12:early-earth:stromatolite}{stromatolieten}
van de Pilbara (3.43); de \omterm{def:g12:early-earth:bif}{gebande ijzerformaties} van de Hamersley Ranges (meer dan
2.5); het einde van de \omterm{def:g12:early-earth:goe}{Grote Oxidatie} (2.06); het einde van de meeste \omterm{def:g12:early-earth:bif}{gebande
ijzerformaties} (1.8).
\end{solution}

\begin{solution}{exo:g12:early-earth:3}
Bij $t = 0$ is $L = L_\odot/(1 + 0.4) = L_\odot/1.4 \approx 0.714\,L_\odot$:
ongeveer \qty{29}{\percent} zwakker, bijna een derde minder fel dan vandaag.
\end{solution}

\begin{solution}{exo:g12:early-earth:4}
Dat de lucht toen bijna geen zuurstof bevatte. Pyriet wordt door zuurstof vernietigd;
zijn korrels konden alleen door rivieren worden meegerold en ongeschonden worden afgezet als
de lucht en het water aan het oppervlak er geen hadden.
\end{solution}

\begin{solution}{exo:g12:early-earth:5}
Waterdamp, die condenseerde tot de oceaan, en koolstofdioxide, dat
in veel grotere hoeveelheden dan vandaag in de lucht bleef en de Aarde door
zijn broeikaseffect verwarmde.
\end{solution}

\begin{solution}{exo:g12:early-earth:6}
3 miljard jaar geleden: $1 - t/t_\odot = 3/4.567 \approx 0.657$, $L \approx
1/1.263 \approx 0.792\,L_\odot$, $T \approx 255 \times 0.792^{1/4} \approx
\qty{241}{K}$. 2.45 miljard jaar geleden: $1 - t/t_\odot \approx 0.536$, $L
\approx 1/1.215 \approx 0.823\,L_\odot$, $T \approx 255 \times 0.823^{1/4}
\approx \qty{243}{K}$.
\end{solution}

\begin{solution}{exo:g12:early-earth:7}
\qty{1}{kg} ijzer is $1000/55.85 \approx 17.9$ mol; één zuurstofmolecule
neemt vier ijzerionen, dus $17.9/4 \approx 4.48$ mol
\ce{O2}, dat is $4.48 \times 32.00 \approx \qty{143}{g}$ zuurstof.
\end{solution}

\begin{solution}{exo:g12:early-earth:8}
(1) $2.46 - 2.06 = 0.40$ miljard jaar, 400 miljoen jaar. (2) Van 4.54
tot 2.46 miljard jaar geleden: $2.08/4.54 \approx \qty{46}{\percent}$ van de
geschiedenis van de Aarde. (3) Ongeveer $3.43 - 2.46 \approx 1$ miljard jaar: de zuurstof
die de microben maakten, werd lange tijd verbruikt voordat ze in de
lucht kon blijven.
\end{solution}

\begin{solution}{exo:g12:early-earth:9}
De \omterm{def:g12:early-earth:bif}{gebande ijzerformatie} toont een diepe zee zonder zuurstof, die
opgelost ijzer kon meevoeren, onder een ondiepe verlichte zee met genoeg zuurstof om het
neer te slaan. De rode laag, op het land afgezet, toont zuurstof in de lucht en in het water
aan het oppervlak, genoeg om het ijzer van haar korrels te laten roesten.
\end{solution}

\begin{solution}{exo:g12:early-earth:10}
In Hamelin Pool, Shark Bay, West-Australië, waar het water twee keer zo
zout is als gewoon zeewater. Een mat van microben, vooral cyanobacteriën, vangt en
bindt sedimentkorrels en laat carbonaat groeien; er groeit een nieuwe mat bovenop,
enzovoort, laag na laag.
\end{solution}

\begin{solution}{exo:g12:early-earth:11}
(1) $21 \times 10^{-5} = \qty{0.00021}{\percent}$. (2) $21/10 =
\qty{2.1}{\percent}$.
\end{solution}

\begin{solution}{exo:g12:early-earth:12}
De \omterm{def:g10:habitable-planet:equilibrium-temperature}{evenwichtstemperatuur} blijft dezelfde als $L(1 - A')$ gelijk is aan $1 - A$:
$1 - A' = 0.706/0.741 \approx 0.953$, dus $A' \approx 0.047$. De Aarde zou
minder dan \qty{5}{\percent} van het zonlicht hebben moeten terugkaatsen, bijna als een
zwart lichaam, wat een oceaanplaneet met wolken en wat ijs niet kan. De
andere uitweg is een sterker broeikaseffect.
\end{solution}

\begin{solution}{exo:g12:early-earth:13}
(1) $273 - 237 = \qty{36}{K}$ om het vriespunt te bereiken; $288 - 237 = \qty{51}{K}$ om
het huidige gemiddelde te bereiken. (2) Het broeikaseffect van vandaag voegt $288 - 255 =
\qty{33}{K}$ toe: de jonge Aarde had een sterker nodig. (3) \qty{0.1}{bar} is een
tiende van de druk van de hele atmosfeer van vandaag, voor koolstofdioxide alleen.
\end{solution}

\begin{solution}{exo:g12:early-earth:14}
(1) Vóór 2450 miljoen jaar geleden bestaat het kenmerk, dus bevatte de lucht minder
dan $10^{-5}$ van de zuurstof van vandaag; na 2090 miljoen jaar ontbreekt het, dus
bevatte de lucht meer. (2) Ja: de \omterm{def:g12:early-earth:goe}{Grote Oxidatie} loopt van ongeveer 2460 tot
2060 miljoen jaar geleden, hetzelfde tijdvak.
\end{solution}

\begin{solution}{exo:g12:early-earth:15}
(1) Ongeveer $4567 - 60 = 4507$ miljoen jaar geleden. (2) Hoogstens $4507 - 4404 =
103$ miljoen jaar na de inslag: de Aarde had een korst en water binnen
ongeveer honderd miljoen jaar na de inslag.
\end{solution}

\begin{solution}{pb:g12:early-earth:1}
\textbf{1.} $4.567 - 4.0 = 0.567$ miljard jaar.

\textbf{2.} $1 - 0.567/4.567 \approx 1 - 0.124 = 0.876$.

\textbf{3.} $L = 1/(1 + 0.4 \times 0.876) = 1/1.350 \approx 0.741$ van
die van vandaag.

\textbf{4.} $0.741 \times 1361 \approx \qty{1008}{W/m^2}$.

\textbf{5.} $1 - 0.741 \approx 0.26$: \qty{26}{\percent} zwakker.

\textbf{6.} $T_{\mathrm{eq}} = \left(S(1-A)/(4\sigma)\right)^{1/4}$.

\textbf{7.} Bij vaste $A$ en $\sigma$ is $T_{\mathrm{eq}}$ een constante maal
$S^{1/4}$, dus $T'/T = (S'/S)^{1/4}$.

\textbf{8.} $0.741^{1/4} = \sqrt{\sqrt{0.741}} \approx \sqrt{0.861} \approx
0.928$.

\textbf{9.} $255 \times 0.928 \approx \qty{237}{K}$.

\textbf{10.} $237 - 273 \approx \qty{-36}{\celsius}$.

\textbf{11.} Dat er 4.4 miljard jaar geleden vloeibaar water aan het oppervlak was, een
oceaan.

\textbf{12.} In ondiep water, zoals levende \omterm{def:g12:early-earth:stromatolite}{stromatolieten} vandaag.

\textbf{13.} De \omterm{def:g12:early-earth:bif}{gebande ijzerformaties} van de Hamersley Ranges, afgezet
op een zeebodem.

\textbf{14.} Ja: \qty{237}{K} ligt ver onder \qty{273}{K}; met de atmosfeer van vandaag
zou het gemiddelde oppervlak bevroren zijn geweest.

\textbf{15.} De gesteenten tonen vloeibaar water onder een Zon die te zwak was om het
vloeibaar te houden met de atmosfeer van vandaag: het \omterm{def:g12:early-earth:faint-young-sun}{probleem van de zwakke jonge Zon}.

\textbf{16.} $288 - 255 = \qty{33}{K}$.

\textbf{17.} Minstens $273 - 237 = \qty{36}{K}$, meer dan vandaag.

\textbf{18.} $A' = 1 - 0.706/0.741 \approx 0.047$: een bijna zwarte planeet,
niet aannemelijk.

\textbf{19.} Koolstofdioxide, veel overvloediger dan vandaag; methaan, gemaakt door
vroege microben, kan hebben geholpen.

\textbf{20.} De \omterm{def:g10:habitable-planet:equilibrium-temperature}{evenwichtstemperatuur} van de Aarde vier miljard jaar geleden was
\textbf{ongeveer 237 K}, zo'n \qty{36}{\celsius} onder het vriespunt.
\end{solution}
